% Transcription — fonds Alexandre Grothendieck, shelfmark 154, batch 1 (pages 1-20).
% Established from the facsimile archives/batches/154/batch-01.pdf, cut from a
% copy of 154.pdf fetched through the project's own relay
% (https://grothendieck-relay.onrender.com/source/154.pdf); PDF page k =
% archivists' page k, no cover sheet. Per the skill transcribe-grothendieck.
% The folder is 175 pages in nine batches (1-20, ..., 161-175), transcribed in
% parallel; this is the first, and no neighbouring file was relied on.
% Notation follows the transcription of folder 80 (Systèmes de pseudo-droites).
%
% Pass: Opus 5.5 (claude-opus-5-5), 2026-09-26 — first pass, unchecked against the pages by a human.
%
% Dating and title. \dating is the folder's own date, copied verbatim from
% src/content/catalogue.ts; the folder belongs to the inventory group
% « Espaces stratifiés » (150-151), but since it has a date of its own the
% group's range is not added. Title from the catalogue, trailing full stop
% dropped; « [Système de pseudo-droites] » is bracketed there because the
% archivists supplied it.
%
% Paper. Pages 1-4 are ink on white paper, written in portrait; pages 6-20
% are fan-fold computer-listing paper written in landscape, the printout
% showing through.
%
% Pages skipped: 5 and 7, computer listings (a job banner of 16 August 1982,
% a JCL header) with nothing in his hand; 10, a blank cream separator sheet.
% Pages summarised: none. Page given only as a figure described in a French
% \note: 20 (a coloured drawing on a listing, no text). Pages transcribed
% from his hand: 1-4, 6, 8, 9, 11-19 (with their figures described in notes).
%
% His pagination and dates. Two runs, each dated at its head:
%   « 18.1.84 » (underlined) at the head of p. 1: run 1-6 on pp. 1, 2, 3, 4,
%   6, 8 (numbered 1, 2, 3, 4, 5, 6), plus p. 9 numbered 6' (a complement);
%   « 23.12.83 » at the head of p. 11: run 1-10 on pp. 11-20 (p. 20, the
%   drawing, carries no visible number, the run otherwise reaching 9 on
%   p. 19). The second run is the earlier in date. All recorded in \note{}s.
%   Numbered items inside the text: the « Quatre possibilités » 1)-4) of
%   pp. 1-6 (and « 5 variantes » X_0, X_1, X'_1, X_2, X_3 on p. 9); a) b) on
%   p. 2; a) b) on pp. 11-12; 1) 2) on pp. 18-19.
%
% What the batch is about.
%   1-9   (18.1.84) « Conventions particulières et constructions relatives
%         aux pos. rel. supersingulières »: the « 5 variantes de la surface
%         de déploiement d'un syst. de pseudo-droites ». X_0 excludes the
%         supersingular (« supercritique ») relative positions, which become
%         vertices of order 4n; X_1 admits them as faces F_i glued along
%         2(n-1) specialisation arcs of combinatorial length 1, with a count
%         of « bords libres » depending on the parity of n (6(n-1) for n
%         even, 2(n-1) for n odd); X'_1 cuts out a small supercritical disc,
%         vertices of order 3 on it; the common defect: the local system of
%         the polygons Pol(D) is defined only off the supercritical faces.
%         X_2 (p. 4, 6) blows up the supercritical vertices, all vertices of
%         order 4, the Pol(D) form a local system, but X_2 does not map to X;
%         a boxed doubt (p. 6) why there is no local system on X_0. X_3
%         (p. 8) replaces D_i by a Möbius band, maps to X, but loses the
%         internal combinatorial interpretation. P. 9: the principal
%         D_2n-covers X~_i attached to the Pol(D) are the « déploiements
%         ultimes »; X'_1, X_2, X_3 the most useful.
%   11-20 (23.12.83) « Structure sur l'ens. des positions relatives
%         P = PR(Σ) d'un syst. Σ = (X, Σ) de pseudodroites »: specialisation
%         arrows (D, T), T a specialisation triangle; injectivity of
%         P~ -> P x P, giving an oriented graph Γ(Σ); Sp_0(D), Gn_0(D) and
%         the polygonal structure on Sp_0(D) ⊔ Gn_0(D) (at least 3 elements;
%         cases D stable / sub-stable; a counterexample with St_4, p. 13);
%         supercritical positions Δ ∈ Σ and specialisation « couloirs »
%         (pp. 13-15); the orientation sets ω(D) and bijections along
%         specialisations; the cellular embedding of Γ(Σ) in a surface
%         Str(Σ), weights of edges and corners summing to 2n at each vertex
%         (p. 16); the polygon Pol^(D) and Ĩ ≃ its edges (p. 17); for
%         Δ = Δ_{i_0}, I \ {i_0} ≃ vertices of P^(Δ_{i_0}), ν + ν_0 = n,
%         and a canonical map from the non-supercritical vertices of Γ(Σ) to
%         antipody-compatible polygonal structures on Ĩ, expected injective
%         (p. 18); comparison of P^(D), P^(D') along a specialisation (p. 19).
%
% Notation fixed once for the batch: \mathcal{X} for his cursive barred X
% (the surfaces X_0 ... X_3), as in folder 80; X, X~ for the plane and its
% double cover; \widetilde{D} for the double cover of a relative position D;
% \underline{\Sigma} for his doubly underlined Σ (the system (X, Σ));
% \operatorname{Pol}(D), \widehat{\operatorname{Pol}}(D), \widehat{P}(D) as
% he writes them; \mathbb{D}_{2n} for his open-faced D (dihedral group), as
% batch 5; « p.r. » for position relative, « ps. dr. » for pseudo-droite;
% \operatorname{Str}(\underline{\Sigma}), \operatorname{Sp}_0, \operatorname{Gn}_0.
% The word read « supercritique » (also « supersingulière », which he
% corrects to it on p. 4) is itself doubtful in places.
%
% Readings to check first: the whole of pp. 3-4 and 6 (fast, many \ill{});
% « faces relatives-ax » on p. 1; the start of p. 8 and « superantiques »;
% the cancelled passage of p. 11; the marginal NBs of pp. 14, 15, 17; the
% middle of the right column of p. 19 (several lines struck through).
\documentclass[11pt,a4paper]{article}
% Le préambule commun à toutes les transcriptions du fonds Grothendieck.
%
% Il tient en une idée : **l'incertitude se compose, elle ne se résout pas.**
% Une transcription qui lisse un mot douteux a détruit la seule chose pour
% laquelle on la faisait — un lecteur doit pouvoir savoir, à chaque endroit,
% ce qui était sur le feuillet et ce qui a été ajouté. Les six macros qui
% suivent sont donc l'appareil critique, et non de la décoration.
%
% Les mêmes commandes sont reconnues par `scripts/render.mjs`, qui produit la
% vue de lecture du site. Une macro ajoutée ici doit l'être aussi là-bas,
% faute de quoi le rendu échouera bruyamment — ce qui est voulu : un rendu
% silencieusement faux est pire qu'un rendu absent.

\ProvidesPackage{grothendieck}[2026/08/08 Transcriptions du fonds Grothendieck]

\usepackage{amsmath,amssymb,amsthm}
\usepackage{mathrsfs}
\usepackage{stmaryrd}
% Les diagrammes commutatifs sont partout dans ce fonds : c'est la langue
% même de la géométrie algébrique de Grothendieck, et une transcription qui
% les rendrait en prose ne transcrirait rien.
\usepackage{tikz-cd}
% Français par défaut — c'est la langue des feuillets — mais l'anglais est
% chargé aussi : la traduction et le résumé sont des documents anglais, et la
% typographie française (espaces avant les deux-points) y serait une faute.
% Ces documents basculent par \selectlanguage{english} en tête de corps.
\usepackage[english,french]{babel}
\usepackage{fontspec}
\usepackage{geometry}
\geometry{a4paper,margin=25mm,marginparwidth=28mm}
\usepackage{marginnote}
\usepackage{xcolor}
% ulem, not soul. `soul`'s \st and \ul are built on a character-by-character
% scan that XeLaTeX's font machinery breaks: the document fails with an
% undefined control sequence rather than degrading. ulem does the same two jobs
% and is Unicode-engine safe. `normalem` stops it hijacking \emph, which the
% transcriptions use for Grothendieck's own emphasis.
\usepackage[normalem]{ulem}
\usepackage{hyperref}
\hypersetup{colorlinks=true,linkcolor=black,urlcolor=black,pdfborder={0 0 0}}

% --- Métadonnées du lot -----------------------------------------------------
% Reprises telles quelles par le script de rendu, qui les affiche en tête de
% la vue de lecture. Elles fixent ce que le fichier prétend couvrir : sans
% elles, une transcription flotte sans point d'attache dans un fonds de seize
% mille feuillets.

\newcommand{\folder}[1]{\gdef\grothFolder{#1}}
\newcommand{\batch}[1]{\gdef\grothBatch{#1}}
\newcommand{\leaves}[2]{\gdef\grothFirst{#1}\gdef\grothLast{#2}}
\newcommand{\foldertitle}[1]{\gdef\grothFoldertitle{#1}}
\gdef\grothFolder{?}\gdef\grothBatch{?}\gdef\grothFirst{?}\gdef\grothLast{?}\gdef\grothFoldertitle{}

% --- Le résumé --------------------------------------------------------------
%
% Il ouvre la lecture modernisée, sous le titre « Résumé », et il est composé
% en petit. Ce n'est pas un ornement : c'est le seul passage du document qui
% ne soit pas des mathématiques, et un lecteur doit pouvoir le reconnaître
% avant de l'avoir lu — puis le sauter s'il connaît déjà le sujet, ou s'y
% arrêter s'il ne le connaît pas. Le filet qui le suit marque la reprise du
% corps.

\newenvironment{resume}
  {\small\setlength{\parskip}{0.45em}}
  {\par\medskip\hrule\medskip}

% --- Mots-clés ---------------------------------------------------------------
%
% En anglais, en clôture du résumé de la lecture modernisée : le vocabulaire
% moderne sous lequel chercher ce que les feuillets construisent. Le manifeste
% du site (`npm run manifest`) les extrait pour étiqueter la cote entière —
% cette ligne est la source unique des tags, il n'y a pas de fichier à côté.
\newcommand{\keywords}[1]{\par\smallskip\noindent
  {\footnotesize\textsc{Keywords} --- \textit{#1}}\par}

% --- L'appareil critique ----------------------------------------------------

% Le numéro de feuillet, en marge. C'est l'ancre qui permet de régler un
% désaccord sur une formule en regardant *un* feuillet plutôt qu'une chemise
% de six cents. Le site s'en sert aussi pour tourner les pages du fac-similé
% quand on fait défiler la transcription.
\newcommand{\leaf}[1]{%
  \marginnote{\footnotesize\color{gray!70}\textsf{#1}}%
  \label{leaf:#1}\ignorespaces}

% Illisible : on ne devine pas. Un mot inventé qui se lit comme les autres est
% le pire résultat possible.
\newcommand{\ill}{\textcolor{red!60!black}{[\,\ldots\,]}}

% Lecture douteuse : le mot est donné, et signalé comme incertain.
\newcommand{\uncertain}[1]{\uline{#1}}

% Ajout de l'éditeur — une lettre manquante, un mot sous-entendu.
\newcommand{\add}[1]{\textcolor{blue!50!black}{[#1]}}

% Biffé par Grothendieck lui-même. Ce qu'il a rayé fait partie du document :
% une rature dit souvent où la pensée a changé de direction.
\newcommand{\struck}[1]{\sout{#1}}

% Note du transcripteur, sur l'état matériel du feuillet.
\newcommand{\note}[1]{\par\smallskip{\small\color{gray!60!black}\textit{[#1]}}\par\smallskip}

% Note marginale de Grothendieck, distincte du corps du texte.
\newcommand{\marginal}[1]{\par\smallskip\noindent\hspace{1em}%
  {\small\color{gray!60!black}#1}\par\smallskip}

% --- Titre ------------------------------------------------------------------

\AtBeginDocument{%
  \begin{center}
    {\large\bfseries Fonds Alexandre Grothendieck --- cote n\textsuperscript{o}~\grothFolder}\\[2pt]
    {\itshape\grothFoldertitle}\\[4pt]
    {\small feuillets \grothFirst--\grothLast\ (lot \grothBatch)}\\[2pt]
    {\footnotesize\color{gray!60!black}Transcription établie d'après le fac-similé numérisé
     par l'Université de Montpellier.\\
     Les lectures incertaines sont soulignées, les illisibles notées [\,\ldots\,],
     les ajouts entre crochets.}
  \end{center}
  \vspace{1em}}

\endinput


\folder{154}
\batch{1}
\pages{1}{20}
\dating{1983-1984}
\watermark{Édition de démonstration}
\foldertitle{[Système de pseudo-droites] : notes manuscrites (1983-1984, s.d.)}

\begin{document}

\page{1}
\note{en tête, à gauche, la date \emph{18.1.84}, soulignée ; en haut à droite, son numéro \emph{1}, premier d'une suite qui court jusqu'à la page 8 (numérotée 6) et à son complément 6' (page 9)}

\section*{Conventions particulières et constructions relatives aux pos.\ rel.\ supersingulières}
\note{titre de sa main, en tête de la colonne de gauche}

Quatre possibilités :

1) On n'admet pas les p.r.\ supersing.\ parmi les sommets du graphe du dépl.\ $A$, i.e.\ comme définissant des faces de $\mathcal{X}$. On trouve une surface $\mathcal{X}_0$, dans laquelle les « pos.\ rel.\ supersingulières » correspondent \struck{aux} à des « sommets supersinguliers » --- ce sont les seuls sommets qui sont d'ordre $4n$, au lieu d'être d'ordre $4$.

2) \struck{\ill{}} On admet les p.r.\ supersingulières \add{comme sommets}, et les spécialisations directes \uncertain{ci-bas} d'une p.r.\ sous-stable) vers celles-ci comme arêtes (correspondant, pour les p.r.\ sous-stables voisines, à certains \emph{arcs} \emph{de spécialisation} exceptionnels.
\note{« comme sommets » est écrit au-dessus de la ligne, lecture \uncertain{probable} ; la parenthèse ouverte avant « correspondant » n'est pas refermée}

Finalement, il y a intérêt à les choisir aussi courts que possible (alors qu'on disait plus \uncertain{haut} aussi longs que possible), savoir de longueur combinatoire $1$, et placés : à égale distance entre les arcs de transition (en fait, de générisation) les plus proches). Ainsi, autour de la face supersingulière $F_i$ correspondant à une p.r.\ supersingulière donnée, vont se raccorder (\struck{\ill{}} \add{le long de} $2(n-1)$ arcs de longueur combinatoire $1$ chacun) les $2(n-1)$ faces correspondant aux positions relatives sous-stables proches, sans compter les \struck{positions relatives} \add{faces relatives-ax} stables proches, qui sont comprises \struck{\ill{}} \add{chacune} entre deux faces sous-stables \uncertain{adjacentes}. \ill{} tout $F$ \ill{} un sommet \ill{} considération, \ill{}. Le nombre $2(n-1)$ est d'ailleurs \ill{} commun \ill{} courant, seulement dans le cas stable, alors \uncertain{que} le cas général c'est $2n_i$, où $n_i$ est le nombre de sommets sur $D_i$ ($= \Sigma_i$). Sur la surface $\mathcal{X}_1$ obtenue, tous les
\note{la lecture « faces relatives-ax », écrite au-dessus de « positions relatives » biffé, est très \uncertain{douteuse} ; « $F_i$ » est ajouté au-dessus de « correspondant »}

\marginal{Les 5 variantes de la surface de déploiement d'un syst.\ de pseudo-droites}
\note{encadré de sa main en tête de la colonne de droite}

\emph{NB} Il n'y a pas de « bords libres » (\uncertain{marges}) en un sommet supersingulier \ill{} toute \struck{droite} ps.\ dr.\ aboutissant \ill{} deux faces $F, F'$ incidentes : $s$) en $s$, de positions antipodiques $F, F'$, sur la face $F$ (en $s$)
\note{« $F$, $F'$ » est ajouté en indice au-dessus de « faces » et au bout de la ligne}
\note{figure : un faisceau d'une dizaine de droites (au crayon) passant par un point $s$, l'une d'elles, en rouge, allant de $F$ à $F'$ ; à droite la légende $\mathcal{X}_0$ et « ($n=5$) ». Au-dessous, deux disques cerclés de jaune, portant en rouge un arrangement de pseudo-droites avec deux points marqués aux extrémités d'un diamètre et un trait épais au bas du cercle, séparés par le mot « ou » ; puis un croquis : un arc jaune (la face $F_i$, légendée « face supercritique ») portant trois « éventails » de traits ; des flèches désignent les « faces sous-stables » et les « faces stables » qui s'y raccordent ; légende $\mathcal{X}_1$}

\page{2}
\note{numéroté \emph{2} en haut}
sommets, y compris aux \struck{\ill{}} \add{incidents} deux faces supercritiques, sont d'ordre $4$ --- les sommets sur une telle face $F_i$ correspondant aux \add{$2n_i$} arêtes \uncertain{superficielles} \ill{} création du $D_i$. Pour l'étude des « bords libres » créés, supposons pour fixer les idées $\Sigma$ stable, on trouve suivant la parité de $n$ :

a) $n$ pair, alors chaque face sous-stable donne \emph{deux} bords libres placés aux extrémités de l'arête d'incidence, et chaque face stable qui touche $F_i$ donne de plus un bord libre (\ill{}), ce qui fait donc
\[
2\bigl(2(n-1)\bigr) + 2(n-1) = 6(n-1)
\]
bords libres, ce qui fait un peu beaucoup.

b) $n$ impair, alors on trouve que chaque face sous-stable adjacente \uncertain{donne} $1$ bord libre (au milieu de l'arête incidente commune), $0$ sur les faces stables intermédiaires, en tout donc $2(n-1)$ bords libres --- c'est juste ce qu'il faut pour exprimer les \emph{exactement} \ill{} recollements, \ill{} des faces : \ill{} une des \uncertain{complétions} de $(\mathcal{X}, \Sigma)$ « affinisé » par $D_i$.
\note{« exactement », souligné, est ajouté au-dessus de la ligne}

Le cas $n$ impair est satisfaisant, mais pour le cas \uncertain{pair} tel quel, on n'y a \ill{} redondance des bords libres. Une façon de s'en tirer dans tous les cas, c'est de \struck{prendre} \add{partir de} la surface $\mathcal{X}_0$, et de la modifier en remplaçant chaque sommet supercritique $s$ par un petit disque supercritique. Dans le cas $n$ \struck{\ill{}} \uncertain{pair}, chaque face sous-stable incidente à $s$ contribue $1$ bord libre, et chaque face stable incidente $1$ bord, et il n'en est \ill{} ; donc en tout il y a $2(n-1)$ bords libres qui bordent un $F_s$, \ill{} qu'il forment un \struck{\ill{}} \emph{NB} \struck{le polygone combinatoire \ill{} \ill{}} cons.\ \ill{} à celui obtenu \ill{} $D_i$ --- \ill{}
\note{les dernières lignes de la colonne, corrigées et en partie biffées, se lisent mal}

\emph{NB} Pour le recoll.\ de la face supercritique avec les faces sous-stables adjacentes, on doit encore définir sur $\widetilde{D}_i$ (image inverse de $D_i$ dans $\widetilde{\mathcal{X}}$) des \struck{\ill{}} \emph{arcs de générisation} \struck{\ill{}} \ill{} aux sommets de $D_i$, qu'on déduit de la \uncertain{longueur} combinatoire $1$ --- on sait \uncertain{déf.} les arêtes de la subdivision cellulaire de $\widetilde{D}_i$ déduites de celles données par les sommets. Mais évidemment, il vaut mieux prendre $\widetilde{D}_i$ \ill{} les \uncertain{voisins} de $D_i$, et définir l'arc de générisation relatif : une position relative sous-stable voisine de $D_i$ correspond, sur la figure ci-contre, \uncertain{avec} un arc de longueur $1$ équidistant des deux sommets sur $\widetilde{D}_i$ correspondants à $s$ et des côtés où \ill{} trouve la p.r.\ $D$
\note{figure : un disque dessiné en rouge, traversé par un arrangement de pseudo-droites, avec deux points noirs aux extrémités d'un diamètre horizontal et, en pointillé jaune, l'arc supérieur entre eux ; un trait séparateur au crayon descend entre les deux colonnes}

\page{3}
\note{numéroté \emph{3} en haut}
façon particulièrement jolie de les raccorder, de façon à reconstituer \struck{\ill{}} l'affinisation supercritique de $\Sigma$ sur $F_i$. On trouve une surface $\mathcal{X}'_1$, où les sommets \struck{supercritiques} sur les faces supersingulières sont d'ordre $3$, tous les autres sommets étant d'ordre $4$. Il faudrait expliciter les arcs de spécialisation. Cet \ill{} \uncertain{synthétiques}, en générisation sur les $\widetilde{D}_i$), non seulement sur les p.r.\ sous-stables voisines de $D_i$, mais maintenant aussi sur les p.r.\ stables voisines. Cette fois, on a des arcs de spécialisation de longueur combinatoire $k$ --- ce qui amène à introduire la subdivision barycentrique \uncertain{déduite} de la subdivision cellulaire qui a les $\widetilde{D}$ et les $\widetilde{D}_i$. \struck{Cet}
\note{figure : $\mathcal{X}'_1$ pour $n = 5$ --- un petit disque au bord pointillé de rouge, traversé par un arrangement de pseudo-droites rouges, d'où rayonnent en alternance des traits rouges et jaunes}

Le défaut commun \struck{aux} à ces trois surfaces représentations $\mathcal{X}_0, \mathcal{X}_1, \mathcal{X}'_1$ \uncertain{pour} les p.r.\ est que le système local, via transport parallèle, des \struck{\ill{}} polygones combinatoires $\operatorname{Pol}(D)$ n'est défini ici que sur le complémentaire de l'\uncertain{ensemble} des \add{sommets ou} faces supercritiques, qui sont des lieux de \uncertain{ramification} \ill{}. Un autre défaut, sans doute \uncertain{important}, est qu'on ne peut définir \ill{} \ill{} des \ill{} dans $X$, \ill{} \ill{} (\struck{\ill{}} \ill{} dans le cas de $\mathcal{X}'_1$ --- \struck{\ill{}} à regarder de façon \struck{\ill{}} détaillée\ldots). C'est un défaut qu'il peut être
\note{les deux derniers alinéas, d'une écriture rapide, sont en grande partie reconstruits ; la troisième surface nommée est lue $\mathcal{X}'_1$ d'après la phrase suivante, l'indice étant \uncertain{peu net}}

\page{4}
\note{numéroté \emph{4} en haut}
utile d'introduire d'autres variantes.

3) \struck{\ill{}} Surface obtenue comme $\mathcal{X}_0$ en excluant les p.r.\ super\struck{singulières}\add{critiques}, mais en introduisant, par \ill{} des sommets supercritiques, des arêtes supplémentaires du graphe $\Gamma$ des p.r., savoir la « \uncertain{blow-up} » \ill{} dessin \ill{} « p.r.\ supercritique » \add{$D_i$} et \uncertain{ceci} \ill{} par une p.r.\ voisine sous-stable, \ill{} stable. Cela fait donc des arcs de transition supplémentaires, qui \emph{\uncertain{vérifient}} la \ill{} \ill{} deux arcs de spéc.\ supercritiques \ill{} $2^{\circ}$, mais qui ne sont plus des arcs de spécialisation, \ill{} \ill{} \ill{} \ill{} \ill{} \ill{} doit être considéré comme un arc de transition \emph{de poids $2$}, \uncertain{et} non $1$ comme les autres. On trouve, dans la surface cellulaire \ill{} $\mathcal{X}_2$ obtenue, \struck{de} (\uncertain{déployée} par le graphe dual de $\Gamma$) des sommets qui sont tous d'ordre $4$, \ill{} nouveaux \ill{} correspondant aux \struck{\ill{} \ill{}} \add{\ill{}} (\ill{} \ill{}) \ill{} $s$ \struck{\ill{}} de $D_i$, \struck{\ill{} \ill{} \ill{}} i.e.\ aux arêtes de la subdivision barycentrique \struck{des $\mathcal{X}$} \ill{} $2(n-1)$ dans le cas stable. Il y a \ill{} $D_i$ est plongée dans $\mathcal{X}_2$, \ill{} \ill{} \ill{} la \struck{\ill{}} \add{décomposition} cell.\ des faces par les sommets de la subd.\ bar.\ de celle qui a $D_i$ \ill{} correspondant à des sommets de $\mathcal{X}_2$.
\note{« critiques » est écrit au-dessus de « singulières » biffé ; dans tout le lot nous lisons « supercritique » un mot dont la graphie reste \uncertain{douteuse}. La moitié inférieure de la page porte plusieurs lignes biffées et des ajouts interlinéaires entre accolades, qui ne se lisent qu'en partie. L'indice de $\mathcal{X}_2$, à la dernière ligne, pourrait être un $3$}
\note{figure, en haut à droite : un trait horizontal jaune coupé par cinq barres verticales}

\page{6}
\note{numéroté \emph{5} en haut ; la feuille est écrite en largeur}
Cette carte \add{$\mathcal{X}_2$} se déduit de $\mathcal{X}_0$ en faisant « éclater » les sommets supersinguliers \ill{}, ou encore de $\mathcal{X}'_1$ en découpant \add{l'int.} des faces supersingulières et en identifiant des pts antipodiques sur leurs bords. --- ce qui explicite aussi la façon dont les $D_i$ sont plongées dans $\mathcal{X}$. Cette fois-ci, on a une carte où toutes les faces peuvent être considérées comme \struck{d'ordre} \add{\ill{}} (\struck{\ill{}} $4n$)), \ill{} tous les sommets d'ordre pair \uncertain{cet} $=$ d'ordre $4$). \uncertain{Sur} cette carte, les $\operatorname{Pol}(D)$ forment un syst.\ local \struck{mais} --- alors que sur $\mathcal{X}_0$, c'est \uncertain{rompu} (\uncertain{sous} \ill{} de ramification \ill{}) autour de chaque sommet. Quant aux bords libres, on peut \ill{} \ill{} de $D \subset \mathcal{X}$, ils se raccordent de façon satisfaisante encore. Cette carte d'autre part ne s'envoie \emph{pas} dans $X$ (sauf erreur).
\note{la phrase « c'est rompu \ldots{} autour de chaque sommet » est soulignée d'un long trait ondulé ; un trait rouge dans la marge renvoie à l'encadré qui suit}

\marginal{Mais pour quelles raisons, au juste, ne serait-il pas vrai qu'on a un syst.\ local sur $\mathcal{X}_0$ (ce qui est pourtant faux !) Il y a un pt qui en ce moment m'échappe. Je ne suis plus trop dans le coup, et devrait reconsulter des notes des semaines précédentes.}
\note{remarque de sa main dans la moitié droite de la feuille, entourée d'un grand cercle rouge}

4) La dernière \ill{} concevable \uncertain{consiste} \ill{} présent pour la carte $\mathcal{X}_3$, qui \struck{\ill{}} \add{par} ailleurs

\page{8}
\note{numéroté \emph{6} en haut ; la feuille est écrite en largeur, sur deux colonnes}
ressemble \uncertain{davantage} près : $\mathcal{X}_3$ (les surfaces avec \uncertain{jointures} \uncertain{marquées}, \ill{} comme les surfaces \ill{} \ill{} : $\mathcal{X}_0, \mathcal{X}_1, \mathcal{X}'_1$ le sont aussi elles). On l'obtient en \add{\ill{}} \struck{\ill{}} les positions relatives super\add{critiques} \struck{singulières} \struck{\ill{}} parmi les « faces » admises, mais en \add{considérant} \struck{\ill{} \ill{}}, pour une $D_i$ fixée ayant $n_i$ sommets, $2n_i$ positions relatives super\ill{}critiques correspondantes, suivant la position \add{$D'_i$} du pt d'intersection initial $D_i \cap D'_i$, qui doit être sur une des $2n_i$ facettes sur $\widetilde{D}_i$. On trouve encore que tous les sommets sont d'ordre $4$. La \struck{\ill{}} ps.\ droite $D_i \subset \mathcal{X}_2$ \add{bande (de Möbius)} \struck{\ill{}} est ici remplacée par une \ill{} \ill{} $D_i$. Cela revient, à partir de $\mathcal{X}'_1$, \struck{\ill{}} à \add{après} \struck{depuis} avoir découpé les faces supercritiques \uncertain{à} \uncertain{les} \ill{} au lieu d'identifier directement des pts antipodiques de chaque \ill{} \ill{} des « \struck{\ill{}} bord » à « \uncertain{faces} d'\add{bord} » \struck{\ill{}} une « p.r.\ \ill{} », qui délimite\add{nt} à ce bord \uncertain{près} une bande, cette bande de Möbius.
\note{la colonne de gauche est lue en partie ; « de Möbius », entre parenthèses et souligné, est ajouté au-dessus de « bande »}

Cette fois-ci on aurait une application $\mathcal{X}_3 \to X$, telle que l'image inverse de $X_1$ soit ce qu'il faut, \struck{Mais} y compris les $D_i$ de tout à l'heure. Mais on n'a \uncertain{toujours} pas sur les faces \add{superantiques} de $\mathcal{X}_3$, et n'y a pas cette fois-ci de représentation affinisée de $\Sigma = (X, \Sigma)$ comme dans le cas de $\mathcal{X}'_1$, \struck{mais} plutôt une représentation d'une partie pertinente d'un voisinage tubulaire \uncertain{conv.} de $D_i$ dans $X$. D'autre part, les \struck{\ill{}} faces \ill{} \ill{} singulières \uncertain{naissant} par des \ill{} \ill{} les autres, de sorte \struck{\ill{} \ill{} \ill{}} que $\mathcal{X}_3$ ne se prête pas à une interprétation en termes de la combinatoire interne, du syst.\ local des $\operatorname{Pol}(D)$.
\note{« superantiques » (lecture \uncertain{douteuse}) est ajouté au-dessus de « faces » ; au bas de la colonne, souligné : \emph{TSVP}}

\page{9}
\note{numéroté \emph{6'} en haut : complément de la page précédente, écrit sur un feuillet de listing imprimé}
En plus de ces $5$ variantes $\mathcal{X}_0, \mathcal{X}_1, \mathcal{X}'_1, \mathcal{X}_2, \mathcal{X}_3$, il y a \add{pour chacune} \struck{les} revêtements de ces cartes (ramifiés pour les trois premières) principaux de groupe $\mathbb{D}_{2n}$, associés aux $\operatorname{Pol}(D)$. Ce sont ces revêtements $\widetilde{\mathcal{X}}_i$ qui sont les déploiements ultimes.

Les surfaces les plus utiles me semblent ici $\mathcal{X}'_1, \mathcal{X}_2, \mathcal{X}_3$ et leurs revêtements $\widetilde{\mathcal{X}}'_1, \widetilde{\mathcal{X}}_2, \widetilde{\mathcal{X}}_3$. Peut-être peut-on se passer des $\mathcal{X}_2, \mathcal{X}_3$ et leurs tildés, au profit de $\mathcal{X}'_1$ et $\widetilde{\mathcal{X}}'_1$, et tout ce que $\mathcal{X}'_1$ s'envoie bien dans $X$ comme prévu --- je \add{n'en} suis pas \struck{tout} tout à fait convaincu, \struck{\ill{}} \add{\ill{}} des alentours, \struck{\uncertain{quelques}} \ill{} comment se \ill{} \uncertain{passer} au voisinage du bord d'une face supercritique.
\note{au-dessus de « suis » un « n'en » ; une rature de lettre sur le tilde de $\widetilde{\mathcal{X}}'_1$}

\page{11}
\note{en tête, à gauche, la date \emph{23.12.83} ; en haut au milieu, son numéro \emph{1}, début d'une nouvelle suite (pages 11 à 20, numérotées 1 à 10)}

\section*{Structure sur l'ens.\ des positions relatives $P = \mathrm{PR}(\underline{\Sigma})$ d'un syst.\ $\underline{\Sigma} = (X, \Sigma)$ de pseudodroites}
\note{titre de sa main ; le $\Sigma$ souligné deux fois est rendu $\underline{\Sigma}$}

a) Relations de spécialisation $D \to D'$ (ou de générisation en sens inverse). À vrai dire, on a un ensemble $\widetilde{P}$ de « flèches de spécialisation » ou paires $(D, T)$, où $D$ est des $\mathrm{PR}$.\ et $T$ est un « triangle de spécialisation » pour $D$ i.e.\ \struck{\ill{}} \struck{\ill{}} ens.\ de $\nu \geq 1$ \ill{} sur $D$, \struck{non incidents} \struck{à un sommet de $\Sigma$}, \struck{formant} un arc sur $\widetilde{D}$ (ou \uncertain{ens.} des \uncertain{côtés} de $D$, avec sa structure polygonale) et dont les $\nu+1$ sommets \add{sont non-sommets de $\Sigma$, et} \ill{} \uncertain{consécutifs} de $\widetilde{D}$) \struck{correspondant à $\nu+1$} \ill{} correspondant : les $D_i \in \Sigma$ qui se rencontrent en un même point $s$ \emph{d'ordre exactement $\nu+1$}, et sans sommet de $\Sigma$ dans le triangle $s, r, t$ \ill{} \ill{} \ill{}, sauf $s$\ldots{} [On exclut le cas où $\Sigma$ est trivial, qu'il vaut mieux traiter à part.]
\note{le passage « à un sommet de $\Sigma$ \ldots{} sauf $s$ » est encadré et barré de quatre longs traits obliques ; nous le donnons tel qu'il se lit. Figure, à gauche : un faisceau de droites (en orange) issues d'un point $s$ coupant un segment jaune (légendé « \uncertain{critère-rive} ») entre deux points $r$ et $t$, avec trois points marqués entre eux}

On a une application
\[
\widetilde{P} \xrightarrow{\ (\mathrm{source},\,\mathrm{but})\ } P \times P
\]
qui associe à $(D, T)$ le couple $(D, D')$, où $D'$ est « déduit de $D$ par \uncertain{le} glissement-spécialisation de $D$ au-dessus de $T$ ». Cette application est-elle injective ? Quand on connaît $(D, D')$, on connaît aussi $s$ (qui est l'unique sommet de $\Sigma$ qui est dans $D'$, non dans $D$) donc aussi l'ens.\ \struck{des \ill{}} des \add{voisins} \ill{} de $\Sigma \cup \lbrace D \rbrace$ qui sont inclus dans le triangle $T$ --- \uncertain{savoir} les sommets communs à $D$ et aux \add{\ill{}} droites de $\Sigma$ passant par $s$. D'autre part, on voit qu'aucune des \struck{\ill{}} droites \uncertain{faisant} \ill{} que ps.\ dr.\ de \struck{\ill{}} $\Sigma' = \Sigma \cup \lbrace D \rbrace$) n'est un bigone, \struck{\ill{}} (car sinon $\Sigma$ serait trivial de centre $s$) ce qui implique que les arêtes \struck{\ill{}} \add{contenues} dans $T$ qui joignent $s$ aux sommets \struck{\ill{}} \add{envisagés} sur $D$, sont ici uniques. \struck{d'autre p.} On voit de $=$ que sur $D \cap T$, \ill{} n'y a \ill{} un autre sommet que ceux-là, ce qui indique que l'on \ill{} \ill{} \ill{} et ce qui détermine \uncertain{bien} $T$ : $\widetilde{P} \to P \times P$ est bien injectif. On trouve donc un \struck{graphe} \add{\uncertain{objet orienté}} \ill{} $P$ comme ens.\ des sommets, $\widetilde{P}$ comme ens.\ d'arêtes, chaque arête étant orientée \struck{\ill{}} \add{i.e.\ \uncertain{munie}} d'une origine privilégiée).
\marginal{\emph{NB} \uncertain{on a} \uncertain{toujours} $D \neq D'$, on $D' \cap S = (D \cap S) \cup \lbrace s \rbrace$}

Soit, pour $D \in P$, $\operatorname{Sp}_0(D)$ l'ens.\ de ses \add{(immédiates) ou « successeurs »} spécialisations\struck{, soit $\operatorname{Gn}(D)$ \ill{}}, \struck{de ses} qui correspond à l'ens.\ des triangles de spécialisation \struck{formés} pour $D$. Soit de $=$ $\operatorname{Gn}_0(D)$ l'ens.\ de ses générisations (immédiates), qui correspond à l'ens.\ des couples $(s, \widetilde{\xi})$, où $s \in D \cap S$, et $\widetilde{\xi}$ est un \uncertain{côté} de $D$ en $s$ i.e.\ un pt de $\widetilde{D}$ sur $s$.
\note{$S$ désigne ici, semble-t-il, l'ensemble des sommets de $\Sigma$. Le $\widetilde{\xi}$ de la dernière ligne est \uncertain{lu} ainsi}

\page{12}
\note{numéroté \emph{2} en haut. Figure en tête de la colonne de gauche : cinq faisceaux de pseudo-droites (en orange) coupant une droite horizontale (en jaune), sur laquelle un trait épais vert marque, pour chacun, l'arc de $\widetilde{D}$ (ou le point) que définit la spécialisation ou la générisation}
Chaque élément de $\operatorname{Sp}_0(D)$ définit un arc \struck{(\ill{})} sur $\widetilde{D}$, ses arcs étant mutuellement disjoints ; chaque pt de $\operatorname{Gn}_0(D)$ définit un \struck{\ill{}} \add{pt} de $\widetilde{D}$, ces pt étant \uncertain{tous} différents et non \struck{\ill{}} contenus dans les arcs précédents. Il en résulte une \emph{structure polygonale} sur l'ens.\
\[
\operatorname{Sp}_0(D) \amalg \operatorname{Gn}_0(D),
\]
qui \struck{\ill{}} s'identifie ainsi à l'ens.\ des arcs issus du sommet $D$ du graphe $\Gamma(\underline{\Sigma})$. \struck{Ce polygone a au moins} \add{Cet ens.\ a au moins} trois éléments je crois. \struck{Soit} Il y en a au moins quatre bien sûr si $\operatorname{card} S \cap D \geq 2$, donc il suffit de regarder les cas $\operatorname{card} S \cap D \in \lbrace 0, 1 \rbrace$. Dans le cas d'une $D$ « stable » i.e.\ $\operatorname{card} S \cap D = 0$, il faut prouver qu'il y a au moins \struck{deux} trois triangles de spécialisation distincts, dans le cas \struck{de} $\operatorname{card} S \cap D = 1$, \struck{il f.} (cas sous-stable) il faut montrer qu'il y en a un au moins un.

\note{figure en tête de la colonne de droite : un faisceau de droites orange (et grises) issues d'un point $s$, coupant une droite jaune en $r$ et $t$ ; les lettres $A$ et $B$ marquent les deux grands triangles de part et d'autre}
a) Cas $D$ stable. On sait que $D$ forme au moins \struck{\ill{}} trois triangles (i.e.\ il y a au moins $3$ triangles de $\Sigma'$ adjacents à $D$), considérons-en un. Il est immédiat qu'il est contenu dans un unique triangle de spécialisation $T$. Considérons les deux grands triangles \struck{défi.} $A$ et $B$ (cf.\ figure) définis par $T$, \struck{\ill{}} adjacents le long de l'arc extérieur $(t, r)$ \uncertain{à} $D$, et \struck{\ill{}} \add{incidents aux côtés} $(s, r)$ et $(s, t)$. Dans chacun, il y a un triangle au moins de $\widetilde{\Sigma}$, adjacent à $D$. Chacun de ces triangles est contenu dans un \uncertain{seul} \ill{} unique triangle de spécialisation, $T'$ resp.\ $T''$. D'où les trois triangles cherchés $T, T', T''$.
\note{la phrase « Dans chacun, il y a un triangle au moins de $\widetilde{\Sigma}$ » porte bien un tilde, là où les lignes précédentes écrivent $\Sigma'$}

b) Cas $D$ sous-stable. \add{Soit $s$ le sommet de $\Sigma$ qui est sur $D$.} Il y a au moins \emph{trois} triangles \struck{de} de $\Sigma'$ adj.\ à $D$, \struck{\ill{}} \struck{\ill{} dont les \ill{} \ill{}} \add{considérons une générisation $D'$ de $D$,} sur laquelle on « \uncertain{relit} » un arc \uncertain{$bc$} (cf.\ figure), \struck{\ill{}} dans des triangles de
\note{« trois » est souligné deux fois}

\page{13}
\note{numéroté \emph{3} en haut. Figure en tête : trois pseudo-droites orange concourantes en un point $a$ d'une droite jaune ; au-dessus de $a$, en pointillé jaune, la générisation $D'$, sur laquelle un trait vert marque l'arc $bc$}
spécialisation $(a\,b\,c)$ de $D'$. Il y a au moins deux autres triangles de spécialisation \struck{\ill{}} pour $D'$, qui auront une base sur $D'$ \struck{\ill{}} qui rencontre $(b, c)$ au plus en un sommet. \ill{} l'un d'eux ne \uncertain{rencontre} \ill{} \ill{} \ill{}, on a gagné. Sinon, on \ill{} la figure
\note{figure : deux pseudo-droites orange et un faisceau issu d'un point d'une droite jaune}

Il y a effectivement un contre-exemple avec $\Sigma = \mathrm{St}_4$, et $D$ \struck{\ill{}} la position sous-stable qui \uncertain{évite} les deux \uncertain{carrés} incidents à $a$ \struck{et \ill{}} \add{\uncertain{symétriques}} \uncertain{l'un de l'autre}, \ill{} triviaux.
\note{figure : l'arrangement standard de quatre pseudo-droites (orange) avec la droite jaune $D$ ; quatre régions portent le chiffre $4$ (quadrilatères). $\mathrm{St}_4$ est lu d'après la page ; il désigne vraisemblablement le système standard de $4$ pseudo-droites}

Mais je m'aperçois que dans la description de l'ens.\ des sommets $P = \mathrm{Pr}(\Sigma)$ de $\Gamma$, j'ai oublié les cas les plus importants : les p.r.\ \emph{supercritiques}, correspondant aux droites de $\Sigma$ lui-même ! Il faut que je décide donc quelles sont les spécialisations vers un tel $\Delta$. Je crois clair qu'elles se \uncertain{déduisent} \add{(ou du moins des spécialisations \emph{immédiates},)} \ill{} \ill{} : \emph{position d'une position stable} qui contient\ldots{} on doit avoir $\operatorname{card}(S \cap D) = 1$ i.e.\ $D$ sous-stable. Soit $S \cap D = \lbrace a \rbrace$, il faut que $a$ soit un sommet de $\Delta$. D'autre part, il doit y avoir un secteur (biangulaire) global de $X$, défini par $D, \Delta$, qui soit un « \emph{couloir} » et la spécialisation se fait dans ce couloir, qui prend nom de \emph{couloir de spécialisation}.
\marginal{\struck{\ill{}} \uncertain{je} \uncertain{fais} !}
\note{la note marginale, écrite en biais entre les deux colonnes, est en partie biffée}
\note{figure : un arrangement de pseudo-droites orange, avec la droite jaune $D$ et une droite $\Delta$ qui la coupe en $a$ ; le secteur entre $D$ et $\Delta$ est hachuré au crayon, une flèche courbe va de $D$ vers $\Delta$}

Il peut y avoir un couloir de spécialisation sur chaque \uncertain{côté} de $D$ en $\lbrace a \rbrace$. Quoi qu'il en soit, on considère \struck{les} un couloir de spécialisation \add{pour $D$}, comme un « équivalent » \struck{\ill{}} d'un triangle de sp.\ pour $D$.

\page{14}
\note{numéroté \emph{4} en haut}
Quand il y en a un, l'extension de $D \cdot \lbrace a \rbrace$ \uncertain{entre} \ill{} peut \struck{\ill{}} être bord de triangles de spécialisation pour $D$ (il y en a deux sur la figure ci-dessus). L'ennui, c'est que les \struck{\ill{} de} \add{\ill{} les couloirs et} spécialisation ne forment pas, avec les autres spéc.\ et gén.\ immédiates, une structure polygonale \add{naturelle}. \struck{Mais} \struck{si $D$ est orientée on peut} Peut-être serait-il plus instructif de \uncertain{penser} en $\Delta$, et regarder les \struck{\ill{}} positions \ill{} \ill{} \uncertain{immédiates}. Mettons-nous dans le cas où il \add{(immédiates)} y en a \uncertain{plus} \ill{} « générisations » \ill{} vers de $\Delta$
\marginal{mais si plusieurs $\Delta$}
\note{figure : deux faisceaux de pseudo-droites orange, reliés par une droite orange ; de part et d'autre, des bandes jaunes et vertes, en trait plein et en pointillé, suivent la droite et s'écartent aux sommets, deux petites flèches doubles marquant l'écart}

\ill{}, les \emph{stables}, \struck{\ill{}} \add{et les} \ill{} \uncertain{sous-stables}. \ill{} les stables \uncertain{j'en} vois deux \uncertain{par} \add{$2$} orientations, \struck{\ill{}} chaque arête, correspondant aux \uncertain{deux} \ill{} parallèles le long de l'arête.
\marginal{les stables \ill{} \ill{} les gén.\ \emph{immédiates} \ill{}, \ill{} gén.\ \emph{\uncertain{proches}} \ill{} \ill{} \ill{} stables}
\note{la note marginale, à droite de la figure, est d'une écriture serrée et se lit mal}

Pour les sous-stables, j'en vois également deux par chaque sommet, correspondant aux deux orientations de $X$ en ce sommet. Ainsi l'ens.\ des \struck{\ill{}} générisations correspond à l'ens.\ des sommets et des arêtes dans le \uncertain{vrai} polygone topologique, image inverse de $\Delta$ dans $\widetilde{X}$. On fera attention que \uncertain{le} polygone ainsi \struck{\ill{}} canonique \uncertain{isom}.\ \ill{} polygone des \uncertain{positions} \uncertain{proches} \ill{} \uncertain{regardé} pour \ill{} p.r.\ $D$ non supercritique --- le choix d'un isom.\ dépend de\ldots{} du choix d'une orientation de $\Delta$. \uncertain{Aussi} est-il clair \ill{} que si \struck{\ill{}} $\Delta$ est l'ens.\ des sommets \add{immédiats} \add{de $\Delta$}, l'ens.\ des générisations \ill{} $\Delta$, l'ens.\ des générisations \struck{\ill{}} s'identifie aux sommets d'un polygone comb.\ d'ordre $\underline{2\nu + 2\nu = 4\nu}$,$^{*}$ \ill{} d'ordre $2$ de celui défini par $\Delta$. L'ens.\ des orientations de ce polygone est l'ens.\ des orientations de $\Delta$ \ill{} à l'ens.\ des orientations de $\Delta$ lui-même.

\note{renvoi de sa main au bas de la page : « $^{*}$ ou, d'ordre $2\nu$ seulement en se bornant aux gén.\ \emph{immédiates} »}

\page{15}
\note{numéroté \emph{5} en haut. Figure en tête : un arrangement de pseudo-droites orange et une droite jaune $D$ ; un arc en pointillé vert, partant de $D$, longe une des droites orange jusqu'à un sommet, où il tourne}
Revenant à une p.r.\ $D$ sous-stable passant par $a \in S$, on peut considérer que \struck{\ill{}} (si $\widetilde{D}$ désigne l'ens.\ des \uncertain{côtés}) un couloir de spécialisation correspond à un \emph{arc ouvert} de $\widetilde{D}$, d'extrémités les deux sommets $a', a''$ de $\widetilde{D}$ au-dessus de $a$. \struck{En identifiant} \struck{De façon} Ainsi on peut dire que \struck{les générisations,} les spécialisations de $D$ sont décrites par des arcs ouverts de $\widetilde{D} - \lbrace a', a'' \rbrace$, mutuellement disjoints, des \uncertain{contenus} dans l'un \ill{} \ill{} deux composantes $\widetilde{D}_1, \widetilde{D}_2$ de $\widetilde{D} \smallsetminus \lbrace a', a'' \rbrace$. \struck{Il y a} \add{Ainsi}, on voit qu'elles \uncertain{forment} une structure polygonale sur \ill{} $a', a''$, et \ill{} \ill{} \ill{} $\widetilde{D}_1, \widetilde{D}_2$

Il y a au moins un arc de spécialisation, \uncertain{il} est que le nb total des sommets du polygone de spécialisation-générisation est \uncertain{au} moins égal à quatre. \struck{C'est} Ainsi. Dans le graphe \add{\struck{\ill{}} orienté} $\Gamma(\underline{\Sigma})$ (\uncertain{où} on \ill{} \uncertain{maintenant} pour \struck{\ill{}} comme sommets les p.r.\ supercritiques), \add{pour \ill{} sommet $D$}, \struck{\ill{} \ill{} \ill{}} l'ens.\ des \struck{\ill{}} \ill{} qui en sont \add{$(\simeq \operatorname{Sp}_0(D) \amalg \operatorname{Gn}_0(D))$} est de cardinal au moins $3$ (c'est au moins $4$ si $D$ n'est pas stable). \add{(i.e.\ dans ses générisations)} Il est muni d'une structure polygonale \add{naturelle}, \struck{\ill{}} d'où l'ens.\ des orientations et \add{en \ill{}} l'ens.\ des orientations de $D$. Pour \ill{} flèche \add{$D \to D'$} de spécialisation, on trouve une bijection \uncertain{naturelle} entre $\omega(D)$ et $\omega(D')$.
\marginal{\emph{NB} L'ens.\ des générisations \add{$\operatorname{Gn}_0(D)$} \ill{} $\omega(D)$ s'identifie à l'image inverse de $D \cap S$ dans $\widetilde{D}$, \ill{} pour $D$ non supercritique, $\widetilde{D}$ désigne \uncertain{la} \ill{} des \uncertain{côtés}, et pour $D$ supercritique, \ill{} \ill{} \uncertain{image} inverse dans $\widetilde{X}$. Toutes \uncertain{flèches}}
\note{$\omega(D)$ est lu comme l'ensemble des orientations (deux éléments) du polygone attaché à $D$ ; la note \emph{NB}, ouverte par une accolade dans la colonne de droite, continue sur la page suivante}

\page{16}
\note{numéroté \emph{6} en haut ; les trois premières lignes achèvent la note \emph{NB} de la page précédente}
de spécialisation $D \to D'$ définit une flèche $\operatorname{Gn}_0(D) \hookrightarrow \operatorname{Gn}_0(D')$ injective, dont le complémentaire est formé de deux \struck{sommets} éléments.

\emph{\struck{Plongement}} \uncertain{Plongement} cellulaire canonique de $\Gamma(\underline{\Sigma})$ dans une surface \struck{\ill{}} $\operatorname{Str}(\underline{\Sigma})$ (la « surface structurante »),
\begin{itemize}
\item sommets stables : toutes les arêtes qui en \uncertain{sortent} \emph{partent} de $D$ ;
\item sommets sous-stables : il y a exactement \emph{deux} arêtes qui \emph{arrivent} à $D$ ;
\item sommets critiques : au moins quatre arêtes qui arrivent à $D$ ;
\item sommets critiques \struck{et} extérieurs : toutes les arêtes en $D$ \emph{arrivent} à $D$.
\end{itemize}
\note{figures, à côté de la liste : pour chaque type de sommet, une étoile de flèches --- au sommet stable trois flèches sortantes (deux croquis reliés par $\simeq$), au sommet sous-stable deux flèches épaisses entrantes et deux sortantes, au sommet critique quatre flèches épaisses entrantes ; à gauche, une étoile où cinq flèches épaisses entrent et trois fines sortent, une flèche courbe y renvoyant}

\emph{Toute arête} de $\Gamma(\underline{\Sigma})$ est munie d'un \emph{poids} ($\in \mathbf{N}^{*}$) qui \struck{est le nb d'arêtes de spécialisation} \struck{\ill{} est égal au} nb d'arêtes sur $D$ (au sens de $\underline{\Sigma}'$) interceptées par l'arête de spécialisation) --- dans le cas de la \uncertain{figure} ci-contre, correspondant au triangle de spécialisation $T$.
\note{figure : à gauche, un triangle de spécialisation $T$ hachuré entre deux pseudo-droites orange et la droite jaune ; à droite, un faisceau de droites orange au-dessus d'un arc vert, avec une petite flèche verticale. $\mathbf{N}^{*}$ rend son N ajouré suivi d'une étoile}

\struck{Tout} Secteur angulaire local de $\operatorname{Str}(\underline{\Sigma})$ en un sommet $D$ correspond à un arc de $\widetilde{D}$, \struck{correspondant à l'arc} \add{\uncertain{compris}} entre deux arcs de spécialisation, consécutifs, ou deux \struck{arcs de sommets} \add{\ill{}} des \struck{spécialisations} \add{générisations} consécutifs (sur $\widetilde{D}$, sans arc de spécialisation entre), ou les cas mixtes ; le nb d'arêtes sur cet arc (au sens de $\underline{\Sigma}'$) est appelé \emph{poids} du secteur angulaire. Avec ces données, on trouve une \emph{carte cellulaire} \struck{pondérée} \emph{à graphe orienté}, et à arêtes et « coins » \emph{pondérés}. \uncertain{Chaque} \ill{} \uncertain{dessin} \ill{} \ill{} sommet $D$)
\note{« à graphe orienté » est souligné et en partie barré d'un trait}

\struck{\ill{}} Somme des pondérations des \uncertain{arêtes} \add{\uncertain{tombant}} toutes les arêtes en $D$, et des \add{pondérations des} tous les coins en $D$ \struck{\ill{}} est égal à $2n$ (où $n = \operatorname{card} \Sigma$)
\note{ce dernier alinéa est marqué d'un double trait vertical dans la marge}

\page{17}
\note{numéroté \emph{7} en haut}
Plus précisément, considérons le polygone \uncertain{configuration} déduit du polygone combinatoire $\operatorname{Pol}(D)$ de $D$ sur $\operatorname{Str}(\underline{\Sigma})$, en « éclatant » chaque sommet (correspondant à une arête de $\Gamma(\underline{\Sigma})$ en $D$) en autant de sommets que l'indique son poids, et en \uncertain{plaçant} sur chaque arête (correspondant à un coin de $\operatorname{Str}$ en $D$) autant de pts que l'indique son poids. On trouve un polygone $\widehat{\operatorname{Pol}}(D)$, \struck{\ill{}} \add{dont} les sommets correspondent aux arêtes \struck{\ill{}} sur $\widetilde{D}$, visualisé comme un \uncertain{contour} parallèle à $D$, \add{\uncertain{avec} ses \emph{arêtes} correspondant aux sommets découpés par les $\Delta_i \in \Sigma$ sur $\widetilde{D}$}\ldots{} \struck{les sommets} Si $I$ est l'ens.\ \ill{} des $\Delta_i$, $\widetilde{I}$ l'ens.\ d'indices \ill{} par $\Sigma$, \ill{} $\overrightarrow{\Delta}_i$ orientés, on trouve par l'ens.\ des \ill{} \ill{} correspondance bijective canonique
\[
\widetilde{I} \;\underline{\simeq}\; \text{ens.\ des arêtes de } \widehat{P}(D)
\]
\struck{\ill{}} canonique plus haut d'où la relation. \struck{\ill{}} Comme $\widehat{P}(D)$ est donc d'ordre pair, il admet un antipodisme, $\widetilde{I}$ aussi et la bijection est compatible avec les antipodismes.
\marginal{\emph{NB} Le « \ill{} » $\widehat{P}(D)$ gén.\ \ill{} des sommets \ill{} « \ill{} » \ill{} de l'un des $\widehat{P}(D)$ qui \ill{}\ldots{} \ill{} de $P(D)$, et \ill{} par l'antipodisme.}
\note{la note marginale est écrite en biais au bas de la page, entre les deux colonnes, et ne se lit qu'en partie ; $\widehat{P}(D)$ abrège $\widehat{\operatorname{Pol}}(D)$}

\struck{Ainsi \ill{}} \emph{Il faudrait examiner} : pour le cas de $D$ supercritique, soit
\[
D = \Delta = \Delta_{i_0},
\]
auquel cas il n'y a que des arêtes de générisation, issues de $D = \Delta$, qui correspondent \ill{} \uncertain{pts} de $\widetilde{\Delta}$ ($=$ image inverse de $\Delta$ dans $\widetilde{X}$) qui \uncertain{sont} au-dessus des sommets sur $\Delta$. Choisissons une orientation de $\Delta$, qui permet d'identifier $\widetilde{\Delta}$ en \ill{} \ill{} $\Pi$ \ill{} de $\Delta$. Nous allons affecter chaque arête \add{de $\Gamma(\underline{\Sigma})$} issue de $\Delta$, correspondant \ill{} $\in S \cap \Delta$, des « poids » \struck{\ill{}} (définis par \ill{} \ill{}) \struck{\ill{}} \add{\uncertain{égal}} au \emph{distinctes} de $\Delta$ qui \struck{\ill{}} passent en $s$, i.e.\ le nb des \emph{sommets} découpés au voisinage de $s$ par une « parallèle » (un \uncertain{côté}) à $\Delta$. Si nous « éclatons » chaque arête de $\Gamma(\underline{\Sigma})$, \struck{\ill{}} des sommets de
\note{figures : au milieu de la colonne de droite, deux faisceaux orange sur une droite orange $\Delta$, longée de part et d'autre par des bandes jaunes et vertes (traits pleins et pointillés), les « parallèles » ; au bas de la page, un faisceau orange issu de $s$ coupé par une droite jaune parallèle à $\Delta$}

\page{18}
\note{numéroté \emph{8} en haut}
$\operatorname{Pol}(D)$, en autant de \uncertain{fragments} que l'indique son poids, on trouve un nouveau polygone combinatoire \add{$\widehat{P}(\Delta)$}, d'ordre $2(n-1)$, essentiellement celui découpé par les $\Delta_i$ ($i \neq i_0$) sur une \emph{parallèle} à $\Delta$. D'où on trouve une bijection canonique
\[
I \smallsetminus \lbrace i_0 \rbrace \;\simeq\; \text{ens.\ des \emph{sommets} de } \widehat{P}(\Delta_{i_0}),
\]
compatible avec les antipodismes.
\marginal{\emph{NB} Cette bijection dépend du choix d'une orientation de $\Delta_{i_0}$ --- en le changeant, on transforme le précédent composé avec l'antipodisme.}
\note{au-dessus de « $I \smallsetminus \lbrace i_0 \rbrace$ », une accolade renvoie au mot « canonique »}
\note{figures : un faisceau orange coupant une droite jaune, recouvert d'une spirale au crayon qui s'allonge vers la droite (le secteur est hachuré) ; au-dessous, une droite orange $\Delta_{i_0}$ traversée de quatre faisceaux, et une droite jaune $D$ qui la longe d'un côté puis, franchissant le sommet marqué d'un gros point, passe de l'autre côté, la bande entre elles étant hachurée}

On aimerait comparer la pondération modifiée d'une arête en $\Delta$, avec sa pondération ordinaire. On trouve pour la pondération ordinaire $\nu_0$
\begin{align*}
\nu_0 &= \text{nb d'arêtes découpées sur la génér.\ } \Delta \text{ par } \Sigma \text{ le long du couloir} \\
&= \text{nb d'arêtes en } D = \ldots n - \nu,
\end{align*}
d'où
\[
\nu + \nu_0 = n
\]
\note{le dernier membre de la première formule est surchargé : un premier terme, peut-être $\operatorname{card}(\ldots)$, est écrasé sous « $n - \nu$ »}
donc la pondération modifiée résulte de la pondération ordinaire (\uncertain{si} on connaît $n$\ldots) \ill{} \uncertain{mais} on peut déduire de là \uncertain{la} pondération ordinaire par les formules locales
\[
\textstyle\sum \text{toutes les pondérations en } D = 2n
\]
(\uncertain{où} $D$ est un sommet non supercritique, p.ex.\ stable).

\struck{\ill{}} On a ainsi trouvé une application canonique
\[
\underbrace{\text{Sommets de } \Gamma(\underline{\Sigma})}_{(= \mathrm{Pr}(\Sigma))} \smallsetminus \underbrace{\Sigma}_{\text{p.r.\ supercritiques}} \longrightarrow \begin{array}{c} \text{ens.\ des structures polygonales sur } \widetilde{I} \\ \text{compatibles avec l'antipodisme} \end{array}
\]
et il ne devrait pas être difficile de montrer que cette application est l'injection. Il faudrait de plus, pour bien faire, voir :

1) Comment « \uncertain{raccorder} » les $P(\Delta)$, $\widehat{P}(\Delta)$ le long d'une flèche de spécialisation

\page{19}
\note{numéroté \emph{9} en haut}
et

2) Comment interpréter toute la structure de graphe orienté : plongement cellulaire, avec pondérations des arêtes et des coins, en termes de l'interprétation de ses sommets \struck{\ill{}} \add{comme} des str.\ polygonales sur $\widetilde{I}$.

Pour 1) L'examen \add{($D \to D'$ dans le cas « ordinaire » ($D'$ non supercritique))} de la situation montre que $\operatorname{Pol}(D)$ et $\operatorname{Pol}(D')$ i.e.\ l'ens.\ des arêtes sur $\widetilde{\widetilde{D}}$ et sur $\widetilde{\widetilde{D}}'$ (où le $\approx$ indique la « parallèle ») \add{i.e.\ des \ill{} sur $D$ et sur $D'$} \struck{\ill{}} « sont les mêmes », \struck{\ill{}} sauf \struck{\ill{}} au voisinage des deux segments de points antipodiques correspondant \struck{\ill{}} \add{à} l'\uncertain{extrémité} de l'arête de spécialisation issue de $D$ (\uncertain{et} \uncertain{l'antipodique}), dans $D'$ l'inverse. De façon plus précise,
\note{figures : en haut, une droite jaune $D$ et, au-dessus, la spécialisation $D'$ en pointillé, qui monte vers le sommet d'un faisceau orange (deux petites flèches verticales, un segment vert sur $D$) ; au bas de la page, la même situation avec les « parallèles » $\widetilde{\widetilde{D}}$, $\widetilde{\widetilde{D}}'$ tracées en pointillé brun, deux flèches désignant les sommets $\sigma$ et $\sigma'$ ; dans la marge, deux petites étoiles de pseudo-droites en $D$ et en $D'$ reliées par une flèche orange, avec les secteurs $\alpha$, $\beta$, $\alpha'$, $\beta'$ et des points $a$, $a'$}
on trouve un \emph{isomorphisme canonique} entre $\widehat{P}(D)$ et $\widehat{P}(D')$
\note{ce membre de phrase est marqué d'un double trait vertical}
transformant l'un des sommets $\sigma$ de $\widehat{P}(D)$ provenant de $\vec{a}$ en \struck{\ill{}} \add{l'un} \emph{antipodique} \add{$\sigma'$} des sommets de $\widehat{P}(D')$ \struck{\ill{}} correspondant à $\vec{a}$. Une autre façon de \uncertain{dire} les choses est celle-ci : Désignons par $a'$ l'arête de spécialisation « opposée » à $a$ (en $D'$), on trouve un isom.\ \add{des polygones} $\operatorname{Pol}(D)$ avec \struck{$\operatorname{Pol}(D')$, \ill{} polygone $\operatorname{Pol}(D')$, transformant (\ill{} \ill{} \ill{}\ldots{} à un} $\operatorname{Pol}(D')$ un \uncertain{certain} polygone déduit de $\operatorname{Pol}(D')$ \add{\ill{}} en sommet $a$, de sorte que $\hat{a}$ \struck{\ill{}} \struck{Cet isom.\ $D$ \ill{}} aille sur $a'$. \struck{\ill{}} On peut \add{donc} trouver une unique arête $\gamma$ de $D$, telle que le polygone \struck{\ill{}} $\operatorname{Pol}'(D)$ déduit de $\operatorname{Pol}(D)$ en plaçant un point $a''$ dessus, soit isom.\ à $\operatorname{Pol}(D')$ par un isom.\ prolongeant le précédent. Cet isom.\ est compatible de plus avec les pondérations dans le sens suivant : \add{\ill{} \uncertain{segments} des polygones (\uncertain{arêtes} de $\Gamma(\underline{\Sigma})$)} \struck{\ill{} qui se correspondent} \ill{} \uncertain{deux} \ill{} \ill{} \ill{} $=$ pondérations (\ill{} \ill{} \ill{}), et aussi la pondération de $\alpha$)
\marginal{\uncertain{cercle} : plus bas !}
\note{la marge gauche de la colonne de droite porte, en travers, « \uncertain{cercle} ; plus bas ! » à côté d'un long trait vertical ; plusieurs lignes du milieu de la colonne sont biffées}

\page{20}
\note{page sans texte : seulement, en haut, un dessin en couleurs sur la même situation que page 19 --- quelques pseudo-droites orange dont trois concourent en un point marqué, une droite jaune $D$ en dessous, la spécialisation en pointillé jaune qui monte vers ce point, et des traits bleus (pleins et pointillés) parallèles à $D$ de part et d'autre du sommet}

\end{document}
